% =====================================================================
% 4. BACKGROUND                     target: 3 - 3.5 pages
%
% This is the "teach the reader" section. The examiner asked for a report
% a young researcher outside the field can follow. Everything my method
% builds on is explained HERE, so sections 5 and 6 can stay focused on my
% own work. Nothing in this section is my contribution -- say so.
% =====================================================================

\section{Background}
\label{sec:background}

\begin{guide}
Four subsections: (1) time series classification, (2) multiple instance learning,
(3) the \millet{} framework, (4) how an explanation is measured. Rules for this
section:
\begin{itemize}\itemsep2pt
  \item Use ``we'' or the impersonal form -- this is other people's work.
  \item Explain every symbol the first time it appears.
  \item Prefer one small picture over a long paragraph.
  \item Do not describe the datasets in detail here (\Cref{sec:setup} does that).
\end{itemize}
\end{guide}

% ---------------------------------------------------------------------
\subsection{Time series classification}
\label{sec:bg:tsc}

\begin{guide}
Half a page. Define the task, then name the three standard deep backbones that this
work compares against, and say in one line what each does.
\textbf{Questions to answer:} What makes time series different from images or text?
Why do convolutions work on them? What is a \emph{receptive field} and why does it
matter for long series?
\end{guide}

\todo{Definition paragraph.} A univariate time series of length $\nsteps$ is a vector
$\bag = (\inst_1, \dots, \inst_{\nsteps})$, and time series classification (TSC) means
predicting one label $y \in \{1, \dots, \nclz\}$ for the whole series.

\todo{One paragraph on the standard deep baselines: FCN and ResNet
\citep{wang2017fcn}, InceptionTime \citep{fawaz2020inceptiontime}, and the general
picture given by the review of \citet{fawaz2019review}. One sentence each. Say that
all three end with global average pooling, which is exactly the step that destroys
the per-time-step information.}

\todo{One short paragraph on dilated / temporal convolutional networks
\citep{bai2018tcn} and depthwise-separable convolutions
\citep{chollet2017xception, howard2017mobilenets}, because my encoder is built from
them. Keep it to what the reader needs later.}

% ---------------------------------------------------------------------
\subsection{Multiple instance learning}
\label{sec:bg:mil}

\begin{guide}
Half a page. This is the key idea of the whole report, so explain it with a picture
and an everyday example (a key ring: the ring opens the door if \emph{at least one}
key on it fits; you are told the ring works, not which key).
\textbf{Questions to answer:} What is a bag, what is an instance, what is a pooling
function? Why does MIL give a per-instance score for free? Which pooling functions
exist (mean, max, attention) and what are they good at?
\end{guide}

\todo{Definition paragraph.} In multiple instance learning
\citep{dietterich1997mil}, a \emph{bag} carries the label and the \emph{instances}
inside it do not. Applied to a time series, the bag is the whole series and each
time step is one instance.

\todo{One paragraph on attention-based MIL \citep{ilse2018attention}: the network
learns a weight for each instance, and those weights are readable as importance.}

\begin{figure}[t]
  \centering
  \figph[0.85]{The MIL view of a time series: the series is a bag, each time step is
  an instance, the encoder embeds every instance, and the pooling head turns the
  instance scores into one bag prediction plus one importance value per time step.}
  \caption{\todo{Caption explaining bag / instance / pooling in one sentence each.}}
  \label{fig:mil_view}
\end{figure}

% ---------------------------------------------------------------------
\subsection{The \millet{} framework}
\label{sec:bg:millet}

\begin{guide}
About one page. This is the direct starting point of my work, so be precise, and be
equally precise that it is \emph{not} mine. Present the general model first, then the
Conjunctive head, which is the baseline I try to beat everywhere in this report.
\textbf{Questions to answer:} What does \millet{} change compared with a normal
classifier? Why is the explanation ``free'' and always consistent with the
prediction? What are the weak points I later attack?
\end{guide}

\begin{priorwork}
Everything in this subsection is the work of \citet{early2024millet}. I reuse their
formulation, their metric implementation and their data loaders without change; my
own heads in \Cref{sec:pooling} are written in the same notation so that the two can
be compared line by line.
\end{priorwork}

\paragraph{General form.}
An encoder $f_\theta$ maps the series to one embedding per time step,
\begin{equation}
  \emb_j = f_\theta(\bag)_j \in \R^{\dmodel},
  \qquad j = 1, \dots, \nsteps ,
  \label{eq:encoder_generic}
\end{equation}
and a MIL pooling head turns these $\nsteps$ embeddings into a bag prediction
$\bagpred \in \R^{\nclz}$ together with an interpretation
$\interp \in \R^{\nclz \times \nsteps}$, where $\interp_{k,j}$ says how much time step
$j$ pushed the model towards class $k$.

\paragraph{Conjunctive pooling (the baseline).}
The head used by the \millet{} baseline combines one attention gate with one
per-time-step classifier:
\begin{align}
  a_j       &= \sig\!\left(\mathbf{w}_a^\top \tanh(W_a \emb_j)\right) \in [0,1],
  \label{eq:conj_gate}\\
  \instpred_j &= W_c \emb_j + \mathbf{b}_c \in \R^{\nclz},
  \label{eq:conj_inst}\\
  \bagpred^{\,k} &= \frac{1}{\nsteps} \sum_{j=1}^{\nsteps} a_j\, \instpred_j^{\,k},
  \qquad
  \interp_{k,j} = a_j\, \instpred_j^{\,k}.
  \label{eq:conj_bag}
\end{align}
\todo{Two or three sentences in words: the gate $a_j$ says ``how much should I listen
to this time step'', the classifier $\instpred_j$ says ``what does this time step
argue for'', and the bag score is simply their average over time. Because the label is
\emph{built} from $\interp$, the explanation cannot disagree with the prediction.}

\paragraph{Weak points I build on.}
\begin{guide}
This short list is the hinge of the whole report: each weakness here becomes one of
my pooling heads in \Cref{sec:pooling}. Keep the wording of the two lists identical
so the reader sees the one-to-one mapping.
\end{guide}
\begin{enumerate}\itemsep2pt
  \item \todo{The gate $a_j$ is shared by all classes, so it cannot say ``this point
        supports class A but argues against class B''.}
  \item \todo{The gate is unnormalised and divided by $\nsteps$, so the scale of the
        bag score depends on the length of the series.}
  \item \todo{The aggregator is a fixed mean, so a few strong evidence points get
        diluted by hundreds of background points.}
\end{enumerate}

% ---------------------------------------------------------------------
\subsection{Measuring the quality of an explanation}
\label{sec:bg:metrics}

\begin{guide}
Half a page and very important: most readers have never seen AOPCR or NDCG@$n$.
Explain the \emph{idea} first in one sentence, then the formula.
\textbf{Questions to answer:} Why can we not just look at the pictures? What does
AOPCR do to the series and why does that test the explanation? When can NDCG@$n$ be
computed at all (only when per-time-step ground truth exists, i.e.\ \webtraffic{})?
Is a higher AOPCR always better?
\end{guide}

\paragraph{AOPCR (Area Over the Perturbation Curve at Random).}
\todo{Plain-English sentence first: remove the time steps the model called most
important, one by one, and watch how fast the model's own confidence drops; compare
that with removing random time steps.}
\begin{equation}
  \text{AOPCR} \;=\; \todo{write the formula in \millet{}'s notation},
  \label{eq:aopcr}
\end{equation}
\todo{One sentence: a higher value means the ranking really points at the evidence.
One sentence of warning: AOPCR is not normalised, so its scale follows the scale of
the logits -- values are comparable \emph{between my models}, but not against numbers
printed in another paper trained with a different recipe.}

\paragraph{NDCG@$n$.}
\todo{Plain-English sentence: on \webtraffic{} we know exactly which time steps were
modified, so we can check whether the model ranked those time steps at the top.
NDCG \citep{jarvelin2002ndcg} is the standard score for ``is the true item near the
top of my ranked list''; 1.0 is perfect.}
\begin{equation}
  \text{NDCG@}n \;=\; \todo{formula},
  \label{eq:ndcg}
\end{equation}

\begin{figure}[h]
  \centering
  \figph[0.85]{How AOPCR is computed: the perturbation curve of the model's own
  confidence as the highest-ranked time steps are masked, next to the curve obtained
  by masking random time steps. The shaded area between the two curves is the score.}
  \caption{\todo{Caption: reading the AOPCR curve.}}
  \label{fig:aopcr_curve}
\end{figure}

% ---------------------------------------------------------------------
\subsection{Related work in one paragraph}
\label{sec:bg:related}

\begin{guide}
Do not write a full related-work chapter -- this is an internship report, not a
survey. One paragraph is enough: post-hoc explainers, MIL for medical images, and
where my work sits between them.
\end{guide}

\todo{One paragraph. Post-hoc explainers \citep{ribeiro2016lime, lundberg2017shap,
sundararajan2017ig} explain a fixed model afterwards and can disagree with it;
\citet{rudin2019stop} argues for models that are interpretable by design. MIL is
heavily used on whole-slide images, where a slide is a bag and a patch is an instance
\citep{ilse2018attention, hashimoto2020msdamil, li2021dsmil}; two of my pooling heads
adapt ideas from that literature to time series. End with one sentence positioning my
work.}
